% §8 Discussion & Limitations — target 0.5 page

\section{Discussion and Limitations}
\label{sec:discussion}

\paragraph{Where LLMs help, where they don't}
\systemname{} surfaces an empirical pattern: LLM agents reach high recall on
misconfigurations, default credentials, and information disclosure — findings
that are textually expressible and well represented in their training corpora.
They consistently underperform on binary industrial protocols (Modbus, CoAP,
LoRaWAN) and on cross-device inferences that require holding an aggregated
network view. This suggests a hybrid deployment model: LLM agents for
configuration auditing, traditional scanners for CVE coverage, and human
operators for the OT/ICS layer.

\paragraph{Implications for agentic systems}
The harness's per-device and per-vulnerability micro-agent structure is one
answer to context scaling, but it cedes cross-device inference. Designs that
keep a global aggregator agent in the loop without paying the full
network-context cost are an open research direction.

\paragraph{Limitations}
(i)~\emph{Simulated infrastructure.} \systemname{} scenarios emulate IoT
devices on Linux VMs; real constrained hardware introduces embedded-specific
attack surfaces (bootloaders, RTOS, firmware) we do not cover.
(ii)~\emph{Evaluator leniency.} Bonus categories are accepted as non-FP even
without ground-truth matches, which inflates precision.
(iii)~\emph{Model variance.} With $k{=}3$ runs per configuration, confidence
intervals on F1 remain wide for some scenarios.
(iv)~\emph{Prompt sensitivity.} We fix a single prompt family across models;
per-model prompt tuning would likely raise scores.
(v)~\emph{Static architectures.} The five patterns are fixed; real deployments
evolve over time and include devices not represented here (RFID badge readers,
industrial PLCs, medical devices).
